%!TEX root = ../thesis.tex
% ******************************* Thesis Appendix A ****************************
\chapter{The Wider Primate Reward Circuit} 

\ifpdf
    \graphicspath{{Appendix1.3/Figs/Raster/}{Appendix1.3/Figs/PDF/}{Appendix1.3/Figs/}}
\else
    \graphicspath{{Appendix1.3/Figs/Vector/}{Appendix1.3/Figs/}}
\fi

In the introduction to this thesis, we provide a brief history of the development of the idea that a representation of utility can be found in the dopaminergic neurons of the midbrain. We focus solely on this subcircuit (and restrict the supporting evidence to rodent, primate, and human studies\footnote{There is evidence for dopamine’s role in guiding reinforcement in fruit flies (\cite{waddellReinforcementSignallingDrosophila2013}) and sea slugs (\cite{reyesReinforcementVitroAnalog2005}) although these are unlikely to be true homologues of the primate reward system. Instead, it is probably likely that conserved dopaminergic systems represent a path of least resistance for the evolution of reward circuits (\cite{barronRolesDopamineRelated2010}).}) in the introduction, for the sake of brevity and because the neuronal data presented in Chapter 4 of this thesis is recorded from these projections. However, economic decisions are universal and complex; it would be outrageous to suggest that such a small population of neurons represents the whole of the neuronal reward circuit. In this appendix, we consider other brain regions of interest for researchers looking for neuronal correlates of reward, focusing on the cortical-basal ganglia circuit.

As illustrated in Figure \ref{fig:da_neuroanatomy}, dopaminergic pathways that originate in the ventral tegmental area (VTA) and the substantia nigra pars compacta (SNc) extend widely to these regions involved in economic decision-making. These pathways form two major circuits: the nigrostriatal projections (shown in red) from the SNc to the striatum and the mesolimbic-mesocortical projections (shown in blue) from the VTA to the prefrontal cortices, the nucleus accumbens and the anterior cingulate cortex.

Dopamine neurons are hypothesised to form a ‘reward retina’ that transmits a signal of how a reward deviates from expectation to these downstream regions of the brain which integrate additional information (\cite{schultz2009midbrain}). Cell populations in the two highlighted pathways receive distinct inputs from a wide range of brain areas, with dopaminergic neurons of the VTA receiving the majority of the input back from the ventral striatum and lateral orbitofrontal cortex, as well as inhibitory input from the lateral habenula (\cite{matsumotoRepresentationNegativeMotivational2009}; \cite{matsumotoTwoTypesDopamine2009}). In contrast, SNc receives most of its input from the somatosensory and motor cortices, as well as the subthalamic nucleus, the dorsal striatum and the nucleus accumbens (\cite{watabeuchidaWholeBrainMappingDirect2012}).

\begin{figure}[p] 
\centering    
\includegraphics[width=1\textwidth]{AppendixC/Figs/neuroanatomy_fig.png}
\caption[The dopaminergic pathways of the ventral tegmental area (VTA) and substantia nigra pars compacta (SNc) project to most regions implicated in economic decision making.]{\textbf{The dopaminergic pathways of the ventral tegmental area (VTA) and substantia nigra pars compacta (SNc) project to most regions implicated in economic decision making.} Dopaminergic nuclei located in the midbrain have wide-ranging projections. The nigrostriatal projections from the substantia nigra pars compacta reach to the striatum (red). The mesolimbic-mesocortical projections originate in the ventral tegmental area are reach to the prefrontal cortices as well as the nucleus accumbens and the anterior cingulate cortex. 
}
\label{fig:da_neuroanatomy}
\end{figure}

\section{The Prefrontal Cortices and Reward}

During the early 1990s, Edmund Rolls showed that single neurons in the orbitofrontal cortex (OFC) of macaques responded to flavour (\cite{rollsGustatoryResponsesSingle1990}) and that, crucially for reward encoding, these responses depended on the satiation of the animal in question (\cite{rollsHungerModulatesResponses1989}). \cite{odohertyAbstractRewardPunishment2001} replicated this finding in human studies using functional magnetic resonance imaging (fMRI), showing separate regions on the orbitofrontal cortex in humans showed blood oxygen level-dependent (BOLD) activity mediated by pleasant and aversive taste stimuli, which was expanded by \cite{andersonDissociatedNeuralRepresentations2003} who showed that OFC fMRI activity was associated with an increase in valence, but not intensity, of an odour stimulus.

In two subsequent fMRI studies of humans performing monetary tasks, Brian Knutson showed that the medial prefrontal cortex (mPFC) shows an increase in the BOLD signal upon receiving an award, but a decrease upon withholding of an expected monetary reward (\cite{knutsonRegionMesialPrefrontal2003}; \cite{knutsonDistributedNeuralRepresentation2005}). Moreover, the BOLD signal incorporated both the magnitude and probability of the monetary rewards, and thus appeared to be computing the expected value of an economic stimulus (see Equation \ref{eq:expected_utility1}; \cite{knutsonDistributedNeuralRepresentation2005}). That single cells in these prefrontal cortices\footnote{We refer to both the medial prefrontal cortex (mPFC) and orbitofrontal cortex (OFC) as prefrontal cortices in this appendix, despite these being anatomically distinct regions. We touch upon why different studies might find activation of one or the other when discussing the replication of BOLD signal activity correlation with bid in human experiments using the BDM at the end of Chapter 2 and so will not here. A brief discussion can also be found in \cite{padoaschioppaOrbitofrontalCortexNeural2017}.} encode economic value was later shown by \cite{padoaschioppaNeuronsOrbitofrontalCortex2006} in a task where macaques made choices between amounts of liquid rewards. This study also dissociated the different activity of different single neurons, showing that some responded to the ‘offer value’ (how much juice or water was offered on a trial), while others responded to the ‘chosen value’ (the amount of the chosen liquid that had greater utility and was chosen by the animal). This multiplexing of the value signal including magnitude, probability, and cost was replicated in the macaque prefrontal cortices by \cite{kennerleyEncodingRewardSpace2009}.

In a similar task in which macaques had to choose between two pictures representing a magnitude of a juice reward, neuronal ensemble recordings from the OFC could be decoded to show the oscillation of value encoding between the two choice offers with response times on a single trial correlating with the time spent by the ensemble in the ‘unchosen state’ (\cite{richDecodingSubjectiveDecisions2016}). Furthermore, low current ($25-50\mu A$ stimulation of macaque OFC cells was sufficient to bias the choices in a similar task, while high current stimulation ($\geq 100 \mu A$) did not bias the choices but increased the variability of the choices (\cite{ballestaValuesEncodedOrbitofrontal2020}).

The rich and heterogeneous responses of frontal lobe neurons to tasks was interesting to us, and we strongly considered recording from this region when developing the project due to the increased sophistication of the BDM when compared to traditional choice paradigms. For example, instead of choosing discretely between two offers, a monkey must instead choose between continuous values of a bid. As a monkey then makes their choice, they pass through the other choices which must be continually considered and find an optimal stopping value to place their bid at.

\section{The Amygdala and Reward}

Many of the pioneering human fMRI studies discussed in the last section also found significant changes in the BOLD signal in the human amygdala during economic value tasks (\cite{breiterFunctionalImagingNeural2001}; \cite{odohertyAbstractRewardPunishment2001}\footnote{Though changes in BOLD signal have been found in both these regions during the same task, it is a simplification to conflate the roles of both these regions as we do to some extent here. Both macaque fMRI and single cell electrophysiology studies have found subtle differences in the neuronal activity of the prefrontal cortices and amygdala in Pavlovian (\cite{saezDistinctRolesAmygdala2017}) and credit-assignment (\cite{chauContrastingRolesOrbitofrontal2015}) tasks.}). The amygdala are a collection of nuclei located in the medial temporal lobe of mammals most commonly implicated in fear processing (\cite{ledouxAmygdala2007}). Once the amygdala was implicated in economic decisions, studies found this fear processing could be extended to uncertainty of a potential loss (\cite{hsuNeuralSystemsResponding2005}; \cite{kahnRoleAmygdalaSignaling2002}), an idea extended by findings that humans with bilateral amygdala lesions show markedly reduced loss aversion (see Appendix A for a discussion of Prospect theory; \cite{demartinoNeurobiologyReferenceDependentValue2009}).

As with the prefrontal cortices, single cell recordings from non-human primates have supported the idea that neurons in the amygdala encode economic value. \cite{bermudezRewardMagnitudeCoding2010} used a Pavlovian task in which the display of stimuli predicted a juice reward of three possible magnitudes, finding that the activity of a population of amygdala neurons increased their activity after the delivery of the juice and that one third of these had changes in activity that correlated with the magnitude of the juice received. In a replication of \cite{padoaschioppaNeuronsOrbitofrontalCortex2006}, \cite{jezziniNeuronalActivityPrimate2020} found that basolateral amygdala (BLA\nomenclature{BLA}{BasoLateral Amygdala}) neurons also showed heterogeneous responses in a choice task, though with many more neurons coding the chosen value than in the OFC.

In Appendix B, we discuss the idea that value arises from progress toward a goal (\cite{juechemsWhereDoesValue2019}). Interestingly, studies in humans and macaques have implicated the amygdala in the progression to some internal stopping point on a ‘spend-save’ task (\cite{grabenhorstPrimateAmygdalaNeurons2016}; \cite{zangemeisterNeuralBasisEconomic2016}). Both humans and monkeys could choose to 'save' (and accumulate liquid rewards with interest) during trials or 'spend' (and consume the saved liquid). In humans, BOLD activation in the amygdala encoded the length and value of save sequences, and single BLA neurons in macaques showed ramping phasic activity as a sequence of save trials progressed toward the monkey's choice to consume their reward.

Perhaps most intriguing for studies involving auction theory is that the amygdala appears to be involved in processing economic and social stimuli (see \cite{rutishauserPrimateAmygdalaSocial2015} for a review). In Appendix A we discuss the ‘winner’s curse’ in auction theory: the idea that in an auction the winner is likely to have overpaid for a good (\cite{capenCompetitiveBiddingHighRisk1971}). The auctions presented in this thesis use a random computer bid, as is typical for the BDM paradigm; however, it is possible to imagine future studies which involve multiple organisms bidding to receive one reward per trial. BLA neurons have been shown to respond to the value of rewards for oneself and to those given to another monkey in a donation game, but not to rewards donated by a computer (termed ‘value-mirroring’ neurons by the authors; \cite{changNeuralMechanismsSocial2015}). Moreover, in a learning task, the amygdala neurons track the value learned from the primates' own choices as well as those made by another monkey in view (\cite{grabenhorstPrimateAmygdalaNeurons2019}. These neurons are then able to simulate the choices made by a different animal, a strikingly similar paradigm to how human’s can optimally bid on auctions by simulating the discounting of other actor’s bids (\cite{wilsonCompetitiveBiddingAsymmetric1967}).

\section{Other Structures Implicated in the Reward Circuit}

We focus on the prefrontal cortices and the amygdala here, as we find the most striking evidence for value encoding in these areas. The entire reward circuit of primates is known to be more complex (\cite{haberRewardCircuitLinking2010}) and involves the striatum (\cite{samejimaRepresentationActionSpecificReward2005}; \cite{straitRewardValueComparison2014}), lateral habenula (\cite{matsumotoRepresentationNegativeMotivational2009}), anterior cingulate cortex (\cite{azabCorrelatesEconomicDecisions2018}), and lateral hypothalamus (\cite{harrisLateralHypothalamicOrexin2007}) among others.

Indeed, as noted in the introduction to this thesis, for an organism in the world economic decisions are both complex (integrating multiple sources of sensory information and contexts) and crucial (to the reproductive fitness of the animal) and so we should expect them to recruit much of the brain's total processing power and therefore multiple brain regions.